\documentclass[11pt]{article}
\usepackage[margin=1in]{geometry}
\usepackage{xcolor}
\usepackage{enumitem}
\usepackage[hidelinks]{hyperref}
\usepackage{parskip}
\usepackage{times}
\setlist[itemize]{leftmargin=1.6em,topsep=2pt,itemsep=1pt}

\newcommand{\reviewer}[1]{\bigskip\noindent{\Large\textbf{Reviewer~#1}}\par\medskip}
\newcommand{\cmt}[2]{\medskip\noindent\textbf{Comment #1.} \textit{#2}\par}
\newcommand{\resp}{\smallskip\par\noindent\textbf{Response:} }
\newcommand{\revs}{\smallskip\par\noindent\underline{\textbf{The specific revisions/updates in the manuscript are detailed below:}}\par}

\begin{document}

\begin{center}
{\Large\textbf{Response to Reviewers}}\\[4pt]
\textbf{Manuscript:} Interpretable Intrusion Detection for IoT: Free
Interpretability, Fragile Open-Set Detection\\[2pt]
\textbf{Manuscript ID:} IoT-69735-2026 \quad\textbullet\quad \textbf{Journal:} IEEE Internet of Things Journal
\end{center}

\noindent\textbf{Dear Editor-in-Chief,}

\smallskip
We thank you and the two reviewers for the careful and detailed reading of our
manuscript. The comments were precise and constructive, and working through them
has improved the rigour and clarity of the paper. Below we respond to every
comment in turn. For each one we restate the concern, describe the change we
made and why it addresses the concern, and give the location of the change in the
manuscript. All added or modified text appears in blue in the accompanying
tracked-changes version, so each revision can be checked directly.

One outcome of this revision is significant enough that we want to state it up
front. Both reviewers asked us to replace our single-seed Concept-Bottleneck
results with a five-seed evaluation (Comments R1.4 and R2.2). When we did, the
``task-coupling OOD penalty'' that had been a headline claim of the original
submission did not reproduce: it rested on one favourable seed. We have removed
the claim rather than defend a result we can no longer support, retitled the
paper, and rebuilt it around the conclusion the five-seed evidence does support,
that representation geometry and scorer choice, not the strength of concept
regularisation, govern open-set detectability. The non-replication is now
reported as a result in its own right. The paper is more honest for the change,
and we thank the reviewers for prompting it.

\smallskip
\noindent\textbf{The principal revisions and updates are summarized as follows:}
\begin{itemize}
  \item Re-evaluated every learned model (MLP, JointCBM $\times4\,\gamma$,
  SequentialCBM, HybridCBM, NeSy-NIDS) over \textbf{five seeds} with mean$\pm$std
  (Table~IV), and replicated the CIC $\gamma$ ablation on a second independent
  five-seed set.
  \item Withdrew the single-seed ``task-coupling'' claim and its figure; retitled
  and re-abstracted the paper around the scorer--representation finding.
  \item Added a \textbf{scorer ablation} (Table~V) comparing per-class
  Mahalanobis with Hamming, Jaccard, energy, and a combined score.
  \item Added \textbf{post-hoc OOD baselines} (MSP, ODIN, energy, deep $k$NN,
  Mahalanobis) on the unconstrained MLP (Table~VI).
  \item Added an \textbf{alternative known/unknown split} on both datasets and a
  \textbf{temporal, session-respecting split} (Table~I).
  \item Replaced the concept-accuracy proxy with a \textbf{direct concept-leakage
  probe} and decoupled concept accuracy from OOD detection (Section~IV-D).
  \item Added a full \textbf{$\lambda_\alpha$ sweep} over six values $\times$ five
  seeds (Table~VII and figure).
  \item Added the \textbf{exact concept and rule definitions with thresholds}
  (Table~II).
  \item \textbf{Quantified} the previously qualitative Fig.~3 and Fig.~4 claims
  (variance--AUROC correlation; threshold-drift reproducibility).
  \item Corrected the computational analysis to \textbf{single-thread, batch-1
  latency} and removed the unsupported on-device claims (Section~IV-E).
  \item Added a research-gap discussion and a ``Scope of the contribution''
  paragraph stating what is and is not novel; redrafted the Abstract,
  Introduction, and Conclusion.
  \item \textbf{Withdrew the ``representation geometry governs detectability''
  reading} after a multi-factor test (new Section~IV-C, Fig.~5). A geometry
  $\times$ architecture $\times$ $\gamma$ factorial (six CBM settings $\times$ five
  seeds $\times$ two datasets) shows that within-class activation variance predicts
  OOD AUROC with \emph{opposite} sign between datasets (positive) and within a
  dataset across architectures ($r{=}-0.84$ on CIC-IoT-2023)---a Simpson's
  paradox---so the earlier two-dataset correlation was a confound. A direct
  annealing manipulation ($k_T$ sweep) moves the variance but AUROC follows only
  weakly and inconsistently. The revised central claim is the negative,
  methodological one: neither concept-fidelity regularisation nor any single
  scorer-independent geometry statistic we examined consistently governs open-set
  detectability, so open-set claims resting on one dataset, seed, or architecture
  can invert when the others are varied.
\end{itemize}

\smallskip
We resubmit the revised manuscript for your consideration and hope the reviewers
find their concerns fully addressed.

\smallskip
\noindent Sincerely,\\
\noindent The Authors

% =====================================================================
\reviewer{1}
\noindent\textbf{Comments to the Author.} \textit{This paper presents a timely
and well-executed study on interpretable open-set intrusion detection for IoT
networks, evaluating two complementary architectures (Concept Bottleneck Models
and Neuro-Symbolic NIDS) with rigorous empirical validation. The work addresses
two critical deployment barriers---lack of interpretability and poor handling of
novel attack families. Recommendation: major revision, particularly regarding
experimental rigor, presentation clarity, and the depth of theoretical
grounding.}

\smallskip
We thank the reviewer for the positive assessment and for comments that sharpened
the paper's rigour and clarity. We address each point below.

\cmt{R1.1}{The paper treats 4 classes as known and 9 as unknown on CTU-IoT-23,
and 5 known/29 unknown on CIC-IoT-2023. The selection criteria for which classes
are known vs.\ unknown is not justified. This choice significantly impacts OOD
difficulty and should be motivated by plausible deployment scenarios.}
\resp The reviewer is right that the partition shapes OOD difficulty and should be
motivated rather than assumed. We have added a ``Known/unknown selection''
paragraph (Section~III-A) that sets out the reasoning. The known set consists of
the coarse traffic archetypes an operator can plausibly label early in a
deployment: a benign baseline together with one representative family from each
dominant malicious category (volumetric DDoS, botnet C\&C/Mirai flooding, and
reconnaissance/scanning). The held-out unknowns are mostly sub-families and
protocol variants of those same categories, which matches the operational
open-set case in which a detector trained on a few labelled exemplars must flag
novel variants of known threat types. We also do not rest on this argument alone:
as explained under Comment~R1.8, the central findings hold under an alternative
known/unknown partition and under a temporal split, so they are not an artefact of
this particular choice.
\revs
\begin{itemize}
  \item Added the ``Known/unknown selection'' paragraph motivating the partition
  by a deployment scenario (Section~III-A).
  \item Cross-referenced the alternative-split and temporal-split checks that
  support the choice empirically (Comment~R1.8; Section~III-A).
\end{itemize}

\cmt{R1.2}{The paper states $\lambda$ was selected by ``coarse grid search over
$\{0.1,0.5,1.0\}$'' (Section III-E). This is insufficient detail---what
validation metric was optimized? How was the grid chosen? The single
$\lambda=0.5$ value used throughout suggests limited hyperparameter exploration.}
\resp The original description was too brief to reproduce, and we have expanded it
in Section~III-E. The weight $\lambda$ scales the concept-supervision term of the
loss; the grid $\{0.1,0.5,1.0\}$ spans an order of magnitude around the operative
regime, and the selection criterion was the known-class validation weighted F1,
the same quantity used for early stopping. Known-class F1 turns out to be almost
invariant to $\lambda$ (within $\pm0.002$ on both datasets), which is why we fixed
$\lambda$ at the middle value and varied $\gamma$ instead. To show this rather
than assert it, we added a $\lambda$-sensitivity table (Table~III, five seeds) for
F1, concept accuracy, and OOD AUROC. F1 stays flat while $\lambda$ trades concept
accuracy against OOD detectability in opposite directions on the two datasets (CTU
AUROC $0.790{\to}0.884$; CIC $0.752{\to}0.724$ as $\lambda$ goes from $0.1$ to
$1.0$). The single value is therefore a considered choice supported by evidence,
not a sign of limited search.
\revs
\begin{itemize}
  \item Clarified the role of $\lambda$, the optimisation metric, and the grid
  rationale (Section~III-E).
  \item Added Table~III reporting F1, concept accuracy, and OOD AUROC across
  $\lambda$ over five seeds.
\end{itemize}

\cmt{R1.3}{The use of Mahalanobis distance on binary rule activations
(Section III-C.4) is theoretically questionable, as the authors note. This
approximation is acknowledged but not justified or evaluated against alternatives
(e.g., Hamming distance, Jaccard similarity).}
\resp Acknowledging the approximation without testing it was a real gap. We now
evaluate it directly with a scorer ablation (Table~V, five seeds) that sets
per-class Mahalanobis against two distances native to the binary simplex (Hamming
and Jaccard to the nearest known-class majority-vote centroid), a logit-space
energy score, and a Mahalanobis+energy combination, all under one protocol. The
outcome is informative. On the well-separated CTU activations every
activation-space scorer performs about the same (Mahalanobis $0.899$, Jaccard
$0.891$, Hamming $0.863$), so the Gaussian assumption is not what limits
performance there. On the near-deterministic CIC activations no activation-space
scorer exceeds $0.64$ and only the logit-space energy score recovers detection
($0.749$). What decides the outcome is the geometry of the representation, not the
metric applied to it. We keep the caveat that Mahalanobis is an approximation on
binary vectors, but the choice is now supported empirically, and we state in the
text that the scorer has to be matched to the representation.
\revs
\begin{itemize}
  \item Added Table~V (scorer ablation: Mahalanobis, Hamming, Jaccard, energy, and
  combined) with discussion (Section~IV-B).
  \item Added an explicit, ablation-backed statement that Mahalanobis is an
  approximation on binary activations (Section~III-C).
\end{itemize}

\cmt{R1.4}{CBMs use ``single seed'' (Section III-E), while NeSy-NIDS uses 5 seeds.
This inconsistency undermines the statistical validity of the CBM results,
particularly for the $\gamma$ ablation study where differences are modest (0.838
to 0.895 AUROC). The authors should have performed multi-seed runs.}
\resp This was the most consequential weakness in the original submission, and we
have corrected it. Every learned model, the MLP, all four JointCBM $\gamma$
variants, SequentialCBM, HybridCBM, and NeSy-NIDS, now reports five-seed
mean$\pm$std (Table~IV). The analysis did more than add error bars. The modest CTU
$\gamma$ differences the reviewer points to ($0.838/0.895/0.883/0.887$) do not
reproduce: over five seeds they become $0.800/0.828/0.835/0.840$ with standard
deviations of $0.023$--$0.044$, and SequentialCBM's earlier CTU value of $0.707$
was an unlucky draw (five-seed mean $0.885$). The CIC ``task-coupling'' effect
that carried part of the original narrative also failed to reproduce; we have
withdrawn it and reframed the paper, as described in the Central Revision below.
\revs
\begin{itemize}
  \item Re-ran every learned model over five seeds; Table~IV reports mean$\pm$std.
  \item Replicated the CIC $\gamma$ ablation on a second independent five-seed set
  and withdrew the task-coupling claim (Central Revision).
\end{itemize}

\cmt{R1.5}{The paper claims ``joint evaluation'' but largely treats the two
architectures independently. The authors should strengthen the novelty narrative
by more explicitly positioning how their integrated evaluation framework provides
unique insights beyond prior work.}
\resp The original framing did not make the value of the joint treatment clear, and
we have reworked it around a unified open-set evaluation framework. Two
structurally different interpretable models, a concept bottleneck and a
neuro-symbolic rule bank, are held to one identical protocol: the same data
splits, OOD scorer, thresholding, and seeds. The point of the design is not to
merge the two architectures but to make their behaviour directly comparable, and
that comparability is what produces the two findings that survive replication.
The dependence of scorer effectiveness on representation geometry is visible only
because the two representations differ under an otherwise identical pipeline, and
the absence of a consistent concept-regularisation effect is visible only because
the effect is tracked across both architectures, both datasets, and multiple
splits. We make this reasoning explicit in the Abstract, the revised contribution
list, and a new ``Scope of the contribution'' paragraph (Section~I), and reinforce
it in the Discussion.
\revs
\begin{itemize}
  \item Added the ``Scope of the contribution'' paragraph and revised the
  contribution list to foreground the unified-protocol insight (Section~I).
  \item Revised the Abstract and Discussion to state what the joint evaluation
  uniquely reveals.
\end{itemize}

\cmt{R1.6}{Can you provide the exact 8 concept definitions for each dataset,
including the threshold values used to generate ground-truth concept labels?}
\resp Yes. Table~II now lists all eight concepts for each dataset together with
the exact threshold expression used to produce every ground-truth label, for
example \texttt{is\_high\_rate} ($\mathrm{Rate}>1000$\,pps) and
\texttt{is\_incomplete\_handshake} ($\mathrm{syn\_count}>2 \wedge
\mathrm{fin\_count}=0$) on CTU-IoT-23. CTU thresholds are given in raw physical
units and CIC thresholds in standardised (z-scored) units, matching the two
preprocessing pipelines, so the concept labels can be regenerated exactly from the
raw features. We also note in the text that, because the labels are deterministic
threshold functions, our reported concept accuracy measures agreement with these
declared thresholds rather than with external semantic ground truth.
\revs
\begin{itemize}
  \item Added Table~II with all eight concept definitions and exact thresholds for
  both datasets, units specified.
  \item Clarified that concept accuracy measures agreement with the declared
  threshold labels (Section~III-B).
\end{itemize}

\cmt{R1.7}{Why was TON\_IoT excluded despite being a widely-used IoT security
benchmark? Would the task-coupling OOD penalty generalize to this dataset?}
\resp The exclusion is a deliberate scoping decision that we should have justified.
Our concept and rule vocabularies are defined over Zeek-style bidirectional flow
statistics (packet and byte counts, inter-arrival times, connection state, and
flag counts), which CTU-IoT-23 and CIC-IoT-2023 share and which allow a controlled
comparison over the same feature interface. TON\_IoT's network features are not
expressed over this interface, so using it would mean re-deriving the entire
concept and rule vocabulary and would confound the same-interface comparison the
study depends on; we have added this reasoning to Section~III-A. On whether the
task-coupling penalty would carry over to TON\_IoT: that claim did not survive
multi-seed replication and has been withdrawn (see the Central Revision), so the
question no longer arises. The mechanisms that do survive, the dependence of OOD
detectability on representation geometry and scorer choice, are properties of the
training objective and representation rather than of any one dataset, and testing
them on further benchmarks is a natural next step.
\revs
\begin{itemize}
  \item Added the TON\_IoT exclusion rationale based on feature-interface
  compatibility (Section~III-A).
  \item Noted that the withdrawn task-coupling claim makes the generalisation
  sub-question moot, and flagged broader-benchmark evaluation as future work.
\end{itemize}

\cmt{R1.8}{How sensitive are the results to the known/unknown class split? Have
you tested alternative splits?}
\resp We tested this two ways. First, we built an alternative CTU known set
(\{Benign, DDoS, Okiru, C\&C\}, promoting C\&C from the unknown pool and demoting
C\&C-HeartBeat) and a matching alternative CIC known set, and re-ran the full CBM
evaluation over five seeds on each. The conclusions hold: every JointCBM variant
stays within $0.005$ of the MLP on known-class F1 ($0.994$ on CTU), and $\gamma$
again fails to buy OOD detection (MLP $0.841$; JointCBM $\gamma{=}0$: $0.721$ down
to $\gamma{=}1.0$: $0.667$), this time in the opposite direction to the main split,
which is itself a sign that the $\gamma$ effect is not consistent. The CIC
alternative split behaves the same way (F1 within $0.002$ of the MLP's $0.819$;
JointCBM AUROC $0.752{\to}0.727$ with $\gamma$). Second, we added a temporal,
session-respecting split (Table~I; see also Comment~R2.3). The qualitative
findings hold across all four dataset/split combinations, which is strong evidence
that they do not depend on the original partition.
\revs
\begin{itemize}
  \item Added alternative CTU and CIC known/unknown splits, each over five seeds
  (Section~III-A).
  \item Added the temporal, session-respecting split (Table~I) as a further check.
\end{itemize}

% =====================================================================
\reviewer{2}
\noindent\textbf{Comments to the Author.} \textit{The manuscript investigates
interpretable open-set IoT intrusion detection using CBMs and a neuro-symbolic
NeSy-NIDS architecture. The topic fits IoT-J well, but the work currently appears
closer to an empirical integration/evaluation of established techniques.}

\smallskip
The reviewer's characterisation is accurate, and it led us to state the nature of
the contribution far more precisely. We address the ten points below.

\cmt{R2.1}{CBM, STE-based differentiable logic, Mahalanobis OOD scoring, and
neural/rule hybridization are all established techniques, while NeSy-NIDS mainly
appears to assemble these components. The authors should explicitly identify
which mathematical or algorithmic mechanism is fundamentally new rather than
application-specific integration.}
\resp We now state the claim of novelty explicitly and keep it modest, in a new
``Scope of the contribution'' paragraph (Section~I). We do not claim that concept
bottlenecks, sigmoid-relaxed threshold rules, Mahalanobis scoring, or neural/rule
hybridisation are individually new, and we say so plainly. The contribution is
empirical and methodological, not a new learning algorithm: a controlled,
multi-seed open-set study that identifies what actually governs
out-of-distribution detection in these models, the coupling between the scoring
rule and the geometry of the representation, and that overturns a plausible but
seed-fragile ``task-coupling'' effect we had ourselves reported. In a literature
where single-seed claims are common, establishing which mechanisms genuinely
transfer to open-set NIDS, under replication and against stronger baselines, is a
useful result, and we have framed the paper around it.
\revs
\begin{itemize}
  \item Added the ``Scope of the contribution'' paragraph delimiting what is and
  is not novel (Section~I).
  \item Revised the Related Work positioning to state that the components are
  established and to place the contribution at the level of the empirical study.
\end{itemize}

\cmt{R2.2}{The CBM experiments use only a single random seed, which seriously
weakens the central conclusions. Differences such as AUROC 0.890 vs 0.895 or
0.883 are too small to interpret without variance estimates, particularly when
NeSy-NIDS itself exhibits $\pm0.016$--$0.028$ AUROC variation.}
\resp This point, shared with Reviewer~1 (Comment~R1.4), drove the central revision
of the paper. All learned models now report five-seed mean$\pm$std (Table~IV), and
the reviewer's specific worry is borne out. With variance estimated, the small CTU
$\gamma$ differences fall well within one standard deviation of each other
(five-seed AUROC $0.800/0.828/0.835/0.840$, std $0.023$--$0.044$), so they cannot
be read as a real ordering, and the CIC ``task-coupling'' peak disappears. We have
withdrawn the affected claims and now report only what is stable across seeds.
\revs
\begin{itemize}
  \item Table~IV now reports five-seed mean$\pm$std for all learned models, and
  claims that do not survive this variance are withdrawn (Central Revision).
\end{itemize}

\cmt{R2.3}{The open-set evaluation protocol may substantially overestimate real
deployment generalization. The manuscript randomly partitions flows before
withholding attack families, but it does not demonstrate separation by device or
capture session.}
\resp Random partitioning can indeed leak session identity between train and test,
and we have added a temporal, session-respecting split to test for it (Table~I). A
strictly device-disjoint attack split is not feasible on CTU-IoT-23, because the
malicious families come from very few source hosts (Okiru, DDoS, and
C\&C-HeartBeat originate from only 4, 7, and 5 source IPs). We instead order each
known class by capture timestamp and train on the earliest 70\% while testing on
the latest 15\%, so test traffic strictly follows training traffic in time. The
analysis also exposed a property of the dataset that we report openly: the CTU
\emph{DDoS} family is temporally bimodal, a high-rate flood capture followed by a
near-zero-rate one, so a timestamp split turns DDoS into a cross-capture task on
which every model, the unconstrained MLP included, scores F1 $0.00$. We therefore
report the three temporally-coherent families, whose weighted F1 drops from
${\approx}0.99$ under a random split to $0.862$ under the temporal one, which
confirms the reviewer's point that random partitioning is optimistic. The central
conclusions still hold: OOD AUROC rises only weakly with $\gamma$
($0.765/0.798/0.831/0.833$ versus MLP $0.795$), the reverse of the alternative-split
ordering, and NeSy-NIDS remains the strongest OOD detector ($0.898$).
\revs
\begin{itemize}
  \item Added the temporal, session-respecting split with per-class F1,
  coherent-family F1, and OOD AUROC (Table~I, Section~III-A), including the DDoS
  bimodality finding.
\end{itemize}

\cmt{R2.4}{The OOD baseline comparison is too limited for the paper's main
open-set claim. Comparing mainly against Mahalanobis-MLP, DT, and RF is
insufficient to demonstrate state-of-the-art open-set capability.}
\resp We have broadened the comparison and corrected the framing of the claim. Five
established post-hoc detectors now run on the unconstrained MLP (Table~VI): maximum
softmax probability, ODIN, energy, deep $k$-nearest-neighbours, and per-class
Mahalanobis. The comparison is telling. Confidence-based scores fail on this task
(inverted on CTU, $0.18$--$0.23$; only modest on CIC, $0.64$--$0.75$), while
feature-space detectors are strong and set the effective ceiling ($k$NN $0.908$ on
CTU / $0.846$ on CIC; Mahalanobis $0.890$ / $0.821$). We do not claim to beat these
black-box detectors, and we now say explicitly that we make no state-of-the-art
open-set claim. What we do show is that interpretable rule and concept spaces keep
competitive OOD separability while remaining auditable, alongside the finding that
the effective scorer is set by the representation.
\revs
\begin{itemize}
  \item Added Table~VI with MSP, ODIN, energy, deep $k$NN, and Mahalanobis on the
  MLP.
  \item Added an explicit statement that no state-of-the-art claim is made, and
  reframed the contribution as competitive-with-auditability (Section~IV-B).
\end{itemize}

\cmt{R2.5}{Mahalanobis modeling over exactly binary NeSy rule vectors is
theoretically questionable and experimentally under-validated. The
eight-dimensional binary activation vectors violate the multivariate Gaussian
assumption, yet the entire NeSy OOD mechanism depends on this approximation.}
\resp A mechanism this central needed explicit validation, which the scorer
ablation from Comment~R1.3 now provides (Table~V). Under one protocol we compare
per-class Mahalanobis against two scorers that make no Gaussian assumption and
work natively on binary vectors (Hamming and Jaccard distance to the nearest class
centroid), a logit-space energy score, and a combination. On the well-separated
CTU activations all activation-space scorers perform about equally, so the Gaussian
assumption is not the limiting factor there; on the near-deterministic CIC
activations none succeeds, whatever assumption it makes. The adequacy of
Mahalanobis is therefore an empirical finding that depends on the representation,
not an unexamined assumption, and the text now presents it that way.
\revs
\begin{itemize}
  \item Table~V validates Mahalanobis against binary-native (Hamming, Jaccard) and
  logit-space scorers; Sections~III-C and IV-B revised accordingly.
\end{itemize}

\cmt{R2.6}{The claim that $\gamma$ regularization eliminates concept leakage is
not demonstrated or unclear how to verify. Table II only reports agreement
between predicted concepts and threshold-generated concept labels.}
\resp Both observations are correct: concept accuracy is only agreement with the
threshold labels and does not measure leakage, and the original claim was not
demonstrated. We replaced the proxy with a direct probe. A linear classifier is
trained to predict the class from (i) the soft concept scores and (ii) their
binarised values, and leakage is the test-accuracy gap, the extra task information
carried by the soft deviations beyond the declared concept values (five seeds).
Measured this way, leakage is negligible on CTU-IoT-23 ($\le0.008$ for every
$\gamma$) and sits at $0.05$--$0.07$ on CIC-IoT-2023 with no reduction as $\gamma$
grows ($\gamma{=}0$: $0.051$; $\gamma{=}1.0$: $0.068$). Since the evidence does not
support the original claim, we have retracted the statement that $\gamma$
eliminates leakage and now present concept accuracy as the property $\gamma$
controls, kept separate from OOD detection.
\revs
\begin{itemize}
  \item Added the soft-versus-binarised linear-probe leakage measurement
  (Section~IV-D).
  \item Retracted the ``$\gamma$ eliminates leakage'' claim and decoupled concept
  accuracy from OOD detectability.
\end{itemize}

\cmt{R2.7}{The claimed $\alpha$-controlled interpretability--OOD Pareto front is
unsupported. Only $\lambda_\alpha=0$ and $0.1$ are evaluated, producing
essentially two operating points, while the manuscript itself admits a complete
sweep is future work.}
\resp Two operating points cannot support a Pareto-front claim, and we have run the
full sweep, $\lambda_\alpha \in \{0,0.05,0.1,0.2,0.5,1.0\}$ over five seeds each
(Table~VII and figure). The results call for a more careful statement. The
regulariser raises the learned gate $\alpha$ (rule reliance) monotonically from
${\sim}0.30$--$0.38$ to ${\sim}0.86$--$0.95$ at no cost in known-class F1, while
OOD AUROC stays flat on CTU (within $0.015$ across the whole sweep) and rises by
about $0.05$ on CIC. We now describe $\lambda_\alpha$ as a control that lets an
operator dial up rule reliance at no accuracy or detection cost, rather than as a
Pareto front, and we removed the earlier suggestion that it could be adjusted
without retraining (each sweep point was trained from scratch), which we mark as
untested.
\revs
\begin{itemize}
  \item Added Table~VII and the $\alpha$-gate sweep figure over six values $\times$
  five seeds (Section~IV-C).
  \item Reworded the claim from a ``Pareto front'' to a ``control knob'' and
  flagged the fine-tuning conjecture as untested.
\end{itemize}

\cmt{R2.8}{100\% rule crispness is not a particularly convincing contribution
because it follows directly from the hard thresholding operation in Eq.~(6).}
\resp Exact binarity follows from the hard threshold and is not, in itself, a
result; we no longer present it as one. Working through this comment also led us to
correct a genuine imprecision in the method description, which we report for
accuracy: the hard $0.5$ threshold is applied only at inference (for the audited
rule vector and OOD scoring), while training uses the annealed sigmoid gates
throughout, and the straight-through gradient is not part of the training path. We
corrected Section~III-C and Algorithm~2 to say this and no longer describe
binarity as a contribution. In its place we report a measure that does carry
information: the fraction of soft activations at $k{=}10$ lying within $0.1$ of
binary before the cut ($0.86$ on CTU, $0.78$ on CIC), which quantifies how close
the audited rule vector is to the evidence the classifier used, and we check that
the thresholding is benign for predictions (the predicted class changes on $0.0\%$
of CTU and $1.2\%$ of CIC validation flows). All previously reported numbers were
produced by this implementation, so none of them change.
\revs
\begin{itemize}
  \item Corrected the method description and Algorithm~2 to state that hard
  thresholding is applied only at inference (Section~III-C).
  \item Replaced the trivial crispness statement with the soft-activation measure
  and a thresholding-benignity check.
\end{itemize}

\cmt{R2.9}{The computational/deployment analysis is too optimistic. The reported
$0.42$--$1.03$\,$\mu$s/sample measurements use an Intel Core i9-9820X desktop-class
CPU with batch size 512, which does not substantiate the statement that the
models fit gateway-class or resource-constrained IoT hardware.}
\resp A batch-512 throughput figure on a desktop CPU says nothing about
gateway-class hardware, and we have removed that claim (Section~IV-E). We now
report two quantities that are both meaningful and honestly scoped: the memory
footprint ($<0.4$\,MB; NeSy-NIDS has $13{,}825$ parameters) and the per-flow,
single-thread, batch-1 latency, which is the conservative regime for a gateway
scoring each flow as it completes (MLP $96$\,$\mu$s, JointCBM $109$\,$\mu$s,
NeSy-NIDS ${\sim}590$\,$\mu$s on the desktop CPU). We state that these are
desktop-class measurements, make no claim about on-device performance, and note
that validation on ARM gateway processors is future work.
\revs
\begin{itemize}
  \item Rewrote Section~IV-E around batch-1 single-thread latency and memory
  footprint.
  \item Removed the gateway-hardware claim and identified on-device ARM
  measurement as future work.
\end{itemize}

\cmt{R2.10}{Fig.~3 attributes CIC's poor NeSy OOD performance to near-deterministic
rule activations and reduced within-class variance, while Fig.~4 claims threshold
movements are semantically consistent; neither claim is supported through
controlled ablations or quantitative correlation analysis.}
\resp Both claims were qualitative, and we have replaced them with quantitative
evidence. For Fig.~3, we measured, across all ten NeSy-NIDS models (five seeds on
each of the two datasets), the relationship between mean within-class
rule-activation variance and OOD AUROC; the two are strongly correlated (Pearson
$r{=}0.976$, $p<10^{-5}$). We note in the text that the two datasets form two
clusters, so the result confirms the between-dataset pattern rather than a
fine-grained law, and we say so explicitly. For Fig.~4, we quantified how
reproducible the learned threshold drifts are across seeds: they are strongly
correlated on CIC (mean pairwise $r{=}0.71$) but seed-dependent on CTU (mean
pairwise $r{=}0.04$). We have tempered the claim accordingly, reading the CIC drift
as systematic and the CTU pattern as indicative only, and no longer assert uniform
semantic consistency.
\revs
\begin{itemize}
  \item Added the within-class-variance versus OOD-AUROC correlation for Fig.~3,
  with an explicit note on the two-cluster caveat (Section~IV-B).
  \item Added the threshold-drift reproducibility correlations and tempered the
  Fig.~4 claim.
\end{itemize}

% =====================================================================
\bigskip
\noindent{\large\textbf{Central Revision (shared across both reviewers)}}\par
\smallskip
\noindent This is the largest change in the revision, and we describe it plainly.
The five-seed evaluation both reviewers asked for (Comments~R1.4 and R2.2) showed
that the ``task-coupling OOD penalty'' reported in the original submission was an
artifact of a single favourable seed. The CIC $\gamma{=}0.5$ peak ($0.727$) rested
on one seed whose $\gamma{=}0$ draw was unusually low ($0.668$; five-seed mean
$0.746$); over five seeds AUROC in fact decreases monotonically with $\gamma$
($0.746/0.739/0.727/0.720$), and an independent second five-seed set reproduces
this with a shallower slope. On CTU the ordering is the reverse, and on the
alternative and temporal CTU splits it reverses again: across four dataset/split
combinations the sign of the $\gamma$ effect flips twice and its magnitude never
exceeds $0.04$. We have removed the task-coupling claim and its figure, retitled
and re-abstracted the paper, and now report the non-replication itself as a
methodological result, that single-seed ablations can manufacture spurious
open-set effects and that open-set claims about interpretable NIDS need multi-seed
replication. The analysis cost us a headline result, but the paper is more honest
and more broadly useful for it, and we thank both reviewers for prompting it.

\end{document}
